\src{192}
\subsection*{\normalfont REMARQUES.}

\noindent 1\textsuperscript{o} Pour les définitions précédant le théorème 6,
il n'était pas nécessaire que $\symup{X}$ soit non singulière, car il n'était pas
indispensable de travailler seulement avec des résolutions \uline{finies}. Mais si $\symup{X}$
est singulière on ne pourra plus affirmer, a priori, que les
$\symup{(\underline{T}_{r}^{s})^{p}(A_1,\ldots;B_1,\ldots)}$ soient des faisceaux
\uline{cohérents}, car dans le complexe de faisceaux
\[
  \symup{\Homsh_{\Osh_X}\bigl(\underline{L}(A_1)\otimes\cdots,\underline{L}(B_1)\otimes\cdots\bigr)}
\]
il y aura une infinité de composantes de degré total donné.

\medskip
\noindent 2\textsuperscript{o} On vérifie facilement que dans les formules
(8.1), on peut remplacer \uline{un} des $\symup{\underline{L}(B_i)}$ par $\symup{B_i}$. Compte tenu de la
proposition 3, cela montre donc qu'on a
\begin{equation}\tag{8.8}
\left\{
\begin{aligned}\symup{\underline{T}_{1}^{1*}(A;B)}&\symup{=\Extsh_{\Osh_X}^{*}(A,B)}\\\symup{T_{1}^{1*}(A;B)}&\symup{=\operatorname{Ext}_{\Osh_X}^{*}(X;A,B)}\end{aligned}
\right.
\end{equation}
En particulier faisant dans (8.6) $\symup{r=s=1}$ et $\symup{A_1=\Osh_X}$, on retrouve le
théorème 3. La formule (8.8) implique aussi
$\symup{T_{0}^{1*}(B)=H^{*}(X,B)}$, $\symup{T_{1}^{0*}(A)=\operatorname{Ext}_{\Osh_X}^{*}(X;A,\Osh_X)}$

\medskip
\noindent 3\textsuperscript{o} On voit sur la formule (8.1) qu'en général les
foncteurs $\symup{\underline{T}_{r}^{s*}}$ et $\symup{T_{r}^{s*}}$ ont des composantes de
degrés positifs \uline{et} négatifs. Utilisant la remarque 2, on voit que si la
dimension de $\symup{X}$ est $\symup{n}$, alors les composantes non nulles de
$\symup{\underline{T}_{r}^{s*}}$ sont comprises entre $\symup{-(s-1)n}$ et $\symup{rn}$, si $\symup{s>0}$,
entre $\symup{0}$ et $\symup{rn}$, si $\symup{s=0}$, et les composantes non nulles de $\symup{T_{r}^{s*}}$
sont comprises entre $\symup{-(s-1)n}$ et $\symup{(r+1)n}$, si $\symup{s>0}$, entre $\symup{0}$ et $\symup{(r+1)n}$,
si $\symup{s=0}$, (et sauf erreur même, si $\symup{r>0}$, entre $\symup{-(s-1)n}$ et $\symup{rn}$, resp. entre
$\symup{0}$ et $\symup{rn}$).

\section*{\centering\normalfont BIBLIOGRAPHIE}

\noindent[1] CARTIER (Pierre). - Les groupes
$\symup{\operatorname{Ext}^{s}(A,B)}$, Séminaire A. Grothendieck : Algèbre homologique,
t. 1, 1957, n\textsuperscript{o} 3.

\noindent[2] GROTHENDIECK (Alexandre). - Sur quelques points d'algèbre
homologique, Tohoku math. J., t. 9, 1957, p. 119-183.

\noindent[3] SERRE (Jean-Pierre). - Faisceaux algébriques cohérents,
Annals of Math., t. 61, 1955, p. 197-278.

\noindent[4] SERRE (Jean-Pierre). - Sur la dimension homologique des anneaux
et des modules noethériens, Proc. Intern. Symp. on alg. number Theory
[1955. Tokyo et Nikko] - Tokyo, Science Council of Japan, 1956 ;
p. 175-189.

\clearpage
